
<<<<<<< HEAD
\chapter{Lista das Características/\emph{Features}\label{chap:Lista-das-Caracter=0000EDsticas}}

Neste capítulo lista-se as características do sistema a ser desenvolvido
(seção \ref{sec:Lista-de-caracter=0000EDsticas}). 

\section{Lista de características <\textcompwordmark <features>\textcompwordmark >\label{sec:Lista-de-caracter=0000EDsticas}}

No final do ciclo de concepção e análise chegamos a uma lista de características
<\textcompwordmark <\emph{features}>\textcompwordmark > que teremos
de implementar.

Após a análises desenvolvidas e considerando o requisito de que este
material deve ter um formato didático, chegamos a seguinte lista:
\begin{itemize}
\item v0.1
\begin{itemize}
\item Ex: O sistema deve plotar gráficos de pressão x temperatura
=======
\chapter{Lista das Caracter�sticas/\emph{Features}\label{chap:Lista-das-Caracter=0000EDsticas}}

Neste cap�tulo lista-se as caracter�sticas do sistema a ser desenvolvido
(se��o \ref{sec:Lista-de-caracter=0000EDsticas}). 

\section{Lista de caracter�sticas <\textcompwordmark <features>\textcompwordmark >\label{sec:Lista-de-caracter=0000EDsticas}}

No final do ciclo de concep��o e an�lise chegamos a uma lista de caracter�sticas
<\textcompwordmark <\emph{features}>\textcompwordmark > que teremos
de implementar.

Ap�s a an�lises desenvolvidas e considerando o requisito de que este
material deve ter um formato did�tico, chegamos a seguinte lista:
\begin{itemize}
\item v0.1
\begin{itemize}
\item Ex: O sistema deve plotar gr�ficos de press�o x temperatura
>>>>>>> 9cc8289ee257736c1a1b3309851f70a5b4b254e3
\item Ex: O sistema deve permitir abrir arquivos num formato especificado.
\item Ex: nome de classe a ser implementada e testada
\end{itemize}
\item v0.3
\begin{itemize}
\item Lista de classes a serem implementadas
\item Testes
\end{itemize}
\item v0.7
\begin{itemize}
\item Lista de classes a serem implementadas
\item Testes
\end{itemize}
\end{itemize}

